\documentclass[a4paper]{article}
\usepackage{titlesec}
\usepackage{geometry}
\geometry{hmargin={2cm,2cm}, vmargin={3.5cm,3.5cm}}
\usepackage{graphicx}
\usepackage[table]{xcolor}
\usepackage{hyperref}
\hypersetup{pdfborder = {0 0 0},
colorlinks,
citecolor=black,
filecolor=black,
linkcolor=black,
urlcolor=blue}
\definecolor{logo}{RGB}{20,99,163}
\usepackage[T1]{fontenc}
\usepackage[utf8]{inputenc}
\usepackage[english,italian]{babel}
\usepackage{fancyhdr}
\usepackage{lastpage}
\usepackage{comment}

% creazione comando subsubsubsection
\titleclass{\subsubsubsection}{straight}[\subsection]

\newcounter{subsubsubsection}[subsubsection]
\renewcommand\thesubsubsubsection{\thesubsubsection.\arabic{subsubsubsection}}
\titleformat{\subsubsubsection}
    {\normalfont\normalsize\bfseries}{\thesubsubsubsection}{1em}{}
    \titlespacing*{\subsubsubsection}
    {0pt}{3.25ex plus 1ex minus .2ex}{1.5ex plus .2ex}

\makeatletter
\def\toclevel@subsubsubsection{4}
\def\l@subsubsubsection{\@dottedtocline{4}{7em}{4em}}
\makeatother

\setcounter{secnumdepth}{4}
\setcounter{tocdepth}{4}
%fine creazione comando


%Comandi di impaginazione uguale per tutti i documenti
\pagestyle{fancy}
\lhead{\includegraphics[scale=0.2]{../../../../Images/logo.png}}
%Titolo del documento
\rhead{\large Verbale Interno del 2021-04-09}
%\rfoot{\thepage}
\cfoot{Pagina \thepage{} di \pageref{LastPage}}
\setlength{\headheight}{45pt}
\renewcommand{\footrulewidth}{0.4pt}
\begin{document}

\begin{titlepage}
    \begin{center}
        \includegraphics{../../../../Images/logo.png}\\
        \vspace{20px}
        \textcolor{logo}{\hrulefill}\\
        \vspace{20px}
        \textbf{\huge\textcolor{logo}{Verbale Interno del 2021-04-09}}\\
        \vspace{10px}
        \textcolor{logo}{\hrulefill}\\
        \vspace{40px}
        \textbf{\Large Informazioni sul documento}\\
        \vspace{20px}
        \begin{tabular}{p{100px} | p{100px}}
            \textbf{Versione}     & 1.0.0                                                                      \\
            \textbf{Responsabile} & Amedeo Meggiolaro                                                          \\
            \textbf{Redattore}    & Samuele De Simone                                                          \\
            \textbf{Verificatore} & Nicolò Giaccone                                                            \\
            \textbf{Destinatari}  & Prof. Tullio Vardanega \newline Prof. Riccardo Cardin, \newline TechSWEave \\
            \textbf{Stato}        & Approvato                                                                  \\
            \textbf{Uso}          & Interno                                                                    \\
        \end{tabular}\\
        \vspace{60px}
        \href{mailto:techsweave@gmail.com}{techsweave@gmail.com}\\

    \end{center}
\end{titlepage}


%INDICE
\newpage
\tableofcontents
\newpage

\section{Informazioni riunione}
\begin{itemize}
    \item Data e ora: 2021-04-09 ore 18:30-19.45.
    \item Luogo: Canale Discord del gruppo tramite videochiamata.
    \item Partecipanti interni: Samuele De Simone, Manuel Varo, Tito Scutari, Nicolò Giaccone, Amedeo Meggiolaro, Marco Barbaresco.
    \item Partecipanti esterni: nessuno.
    \item Segretario: Amedeo Meggiolaro.
\end{itemize}
\section{Argomenti della riunione}
\begin{itemize}
    \item Discussione sul cruscotto: dopo il riscontro del professore riguardo il cruscotto, il gruppo ha discusso su come svilupparlo;
    \item Scelta del servizio di Continuous Integration: il gruppo ha discusso sul miglior servizio da poter utilizzare per la Continuous Integration, scegliendo tra Github Actions, TravisCI, CircleCI e Vercel. Per il backend si è scelto di utilizzare TravisCI, siccome è già stato studiato da alcuni membri, mentre per il frontend si è scelto Vercel, che è un servizio offerto appositamente per piattaforme che utilizzano Next.js.
    \item Discussione sul workflow: per la scrittura del codice si è scelto di utilizzare Github flow.
    \item Discussione sugli incarichi per il weekend: il gruppo, per dividersi il carico di lavoro, ha deciso di suddividere le tecnologie da studiare tra i vari membri, così ognuno sia più specializzato in una o più determinate tecnologie.
\end{itemize}
\section{Resoconto}
\begin{center}
    \textbf{\Large Resoconto riunione}\\
    \vspace{10px}
    \begin{table}[h!]
        \centering
        \rowcolors{2}{logo!10}{logo!40}
        \renewcommand{\arraystretch}{1.8}
        \begin{tabular}{p{200px} p{200px}}
            \rowcolor{logo!70} \textbf{Quesito}           & \textbf{Soluzione}                                                                                                                        \\
            Cruscotto                                     & Il gruppo ha deciso come sviluppare il cruscotto                                                                                          \\                                                                                                                     \\
            Scelta del servizio di Continuous Integration & Il gruppo ha scelto di utilizzare TravisCI per la Continuous Integration del codice per il backend, mentre Vercel per quello del frontend \\
            Discussione sul workflow                      & Si è deciso di utilizzare Github flow                                                                                                     \\
            Discussione sugli incarichi per il weekend    & È stato scelto di suddividere il carico di studio tra i membri.                                                                           \\
        \end{tabular}
    \end{table}
\end{center}

\end{document}